% Ch.9 Discussion (~5-7 pp). Drafted Day 11/12.
% Interprets the four findings of Ch.7-8b, states the limitations promised in
% Ch.5 (HAC vs reverse-regression, EH-MC calibration gap) and Ch.4 (coverage),
% and draws the portfolio / term-premium implications.
\chapter{Discussion}
\label{ch:discussion}

\section{Economic Interpretation}
\label{sec:disc-econ}

\paragraph{What the cycle factor captures internationally.}
The central empirical fact of \Cref{ch:results} is that the trend--cycle
mechanism of \citet{cieslak2015} is not a property of one market but of the
asset class. In every G10 market, the deviation of yields from their
trend-inflation anchor predicts next year's excess return with the same sign
structure---a negative loading on the one-year cycle and a positive loading
on the average cycle---and with loadings that are flat-to-rising across the
maturity curve, the signature of a single premium that prices the whole curve.
The natural reading is the one \citet{cieslak2015} give for the United States,
now with cross-sectional support: trend inflation pins down where investors
expect yields to settle, so the distance of today's curve from that anchor
measures the compensation demanded for holding duration, and that compensation
mean-reverts at business-cycle frequency. The uniformity across eleven
countries with different monetary histories is hard to reconcile with a
data-mined artefact; it is exactly what one expects if the mechanism is
fundamental, as the falling-stars evidence of \citet{bauer2020} implies---once
the slow-moving macro anchors are accounted for, the cyclical deviation is
where the risk premium lives.

The comparison with the forward-rate factor sharpens this interpretation. In
sample, the \citet{cochrane2005} factor fits \emph{better} than the cycle
factor across the G10 ($R^{2}$ of $0.38$ against $0.25$ on average), as a
six-predictor combination of prices should against a two-predictor one. Out
of sample the ranking reverses decisively: of the four candidate factors,
only the \emph{global cycle} factor beats the recursive historical mean
(\Cref{sec:rob-cf-cp}). The macro anchoring is therefore not packaging; it is
load-bearing. Detrending by a slowly-moving, backward-looking inflation trend
produces predictive loadings that are stable enough to estimate in real time,
whereas the forward-rate combination---fit to the fast-moving forward curve
itself---is precisely the kind of in-sample construct that the out-of-sample
test exists to expose. The thesis's evidence thus speaks to the oldest
dispute in this literature, the choice of state vector: between prices
(forwards) and macro-anchored cycles, the cycle wins where it matters, in
real time.

\paragraph{The integration story.}
The second finding---that a single GDP-weighted global cycle factor matches
the explanatory power of eleven national factors and drives the local factor
out in eight of eleven markets---re-establishes the integration result of
\citet{dahlquist2013} on the macro-anchored factor family. The economic
content is that the predictable component of government-bond returns is
compensation for a \emph{common} risk: bond risk premia in the G10 rise and
fall together, and a domestic investor in, say, Sweden is rewarded for
exposure to the global cycle, not to a Swedish one. Two mechanisms plausibly
support this. Economically, the inputs to the factor are themselves
synchronised---inflation trends fell together across the G10 over the sample,
business cycles co-move, and monetary policy responded to common shocks, so
the national cycles share their informative variation. Statistically,
aggregation averages away the estimation noise in eleven separately fitted
generated regressors, which is why the global factor is also the only one
that survives out of sample (\Cref{sec:rob-oos}). These two readings are
complements, not rivals: the GDP-weighted average can only de-noise the
national signals because there is a common signal to recover.

The exceptions are as informative as the rule. The local factor retains
incremental content in Italy, the Netherlands and Belgium, and the Italian
case decomposes cleanly in the subsample analysis: Italian local
predictability is overwhelmingly a 2010--2012 phenomenon, with a horse-race
$t$-statistic above ten and a positive real-time $R^{2}_{\mathrm{oos}}$ of
$+0.21$---the only economically meaningful local out-of-sample result in the
study---while the global factor enters with the wrong sign
(\Cref{sec:rob-italy}). Integration, in other words, is a regime, not a law.
When redenomination and sovereign-credit risk reprice a peripheral market,
the compensation becomes local again, and a GDP-weighted factor dominated by
the core economies cannot price it. This conditionality refines rather than
contradicts \citet{dahlquist2013}: global risk compensation is the rule in
calm integration, and the limit of that rule is sovereign stress.

\paragraph{The role of the FX adjustment.}
The third finding is a contrast with the forward-rate literature and, we
would argue, a substantive one. For \citet{dahlquist2013}, adjusting the
global factor for currency risk matters: their FX-adjusted forward factor
correlates only about $0.50$ with the unadjusted one and recovers
dollar-investor predictability. The cycle-factor analogue built here is
nearly collinear with the unadjusted global factor (correlation $0.99$), so
the FX adjustment, while uniformly helpful in direction, has almost no
independent information to add. A plausible economic explanation lies in
what each factor family retains. The forward-rate factor is built from the
\emph{level} and shape of the forward curve, and cross-country differences in
rate levels are exactly the interest differentials that forecast currency
returns through the failure of uncovered interest parity; an FX adjustment
of a forward factor therefore has carry-relevant information to work with.
The cycle factor, by construction, \emph{removes} the level: detrending by
trend inflation strips out the persistent rate-differential component and
keeps the transitory deviation. What remains forecasts bond premia well but
currencies hardly at all, so projecting dollar returns on global cycles
returns essentially the cycles themselves. On this reading, the
near-collinearity is not a failure of implementation but a property of
macro-anchored factors: the same detrending that gives the cycle factor its
real-time stability in bonds takes away its leverage over currencies. The
regime evidence supports this: the residual dollar-investor predictability
is concentrated before 2008, when rate differentials were wide, and turns
negative in the post-crisis regime of compressed, synchronised policy rates
(\Cref{sec:rob-subsample}).

\paragraph{What ultimately survives.}
Assembling the four findings, the thesis's evidence reduces to one robust
object and a set of conditional ones. The GDP-weighted global cycle factor
predicts international bond returns in sample, out of sample, in every
subsample except the 2008 panic, beats the forward-rate alternative in real
time, and carries economic value for a hedged investor
(\Cref{ch:strategy}). The local factors are real in sample but are, for
nine of eleven markets, not estimable precisely enough to use in real time;
the dollar-investor factors are creatures of the currency regime. The
honest summary is therefore not ``the cycle factor works everywhere'' but
something more specific and more useful: \emph{international bond risk
premia contain a common, macro-anchored, slowly-moving component that is
large enough and stable enough to estimate in real time---provided one
aggregates across markets, measures the inflation anchor cleanly, and does
not ask it to forecast exchange rates.}

\section{Limitations}
\label{sec:disc-limitations}

Five limitations qualify these conclusions; none, we believe, overturns
them, but each bounds how far they should be pushed.

\paragraph{Inference under overlapping returns and persistent regressors.}
All headline $t$-statistics use $18$-lag Newey--West HAC standard errors
\citep{neweywest1987}, whereas \citet{cieslak2015} use the
reverse-regression delta method of \citet{weiwright2013}, which requires the
underlying monthly returns and is not identified from the twelve-month
overlapping panel constructed here (\Cref{sec:meth-inference}). For highly
persistent regressors the two can diverge, and HAC inference at long
horizons is known to over-reject in finite samples \citep{hodrick1992}. The
block-bootstrap intervals \citep{politis1994} and the
expectations-hypothesis Monte Carlo mitigate this---the predictive $R^{2}$
are judged against a simulated no-predictability band, not only against
asymptotic $t$-ratios---but a residual upward bias in the in-sample
significance cannot be excluded. It is precisely because of this that the
thesis leans its conclusions on the out-of-sample evidence, which is immune
to the over-rejection problem: a recursive forecast either beats the
prevailing mean or it does not.

\paragraph{Sample span and a single macro regime.} The panel covers
1989/1994--2026: roughly thirty years of monthly data, but---because returns
overlap---the equivalent of only about thirty non-overlapping annual
observations per country, spanning one secular disinflation, one
zero-lower-bound decade, and one inflation shock. The sample is also
unbalanced (the nine non-US, non-Japanese markets begin only in December
1994), and the five euro-area members share a currency and a monetary
policy from 1999, so the eleven markets contribute far fewer than eleven
independent observations of the ``global'' cycle. The GDP-weighted factor is
to a significant degree a US--euro-area--Japan factor, and the integration
finding should be read accordingly; alternative weightings (equal,
market-capitalisation, principal-component) were not explored.

\paragraph{Construction constraints.} The maturity menu
$\{1,2,4,5,9,10\}$ years is dictated by the available Bloomberg curves, so
individual excess returns exist only for $n\in\{2,5,10\}$ and the
duration-standardisation uses the continuous-compounding convention
$D^{(n)}=n$ (\Cref{sec:meth-deviations}). The FX-adjusted factor is built
top-down on aggregate dollar returns rather than as the projection in the
original proposal. These choices are internally consistent and are held
fixed throughout, but exact level comparisons with the published US results
inherit them.

\paragraph{The EH Monte Carlo calibration gap.} The simulated
expectations-hypothesis benchmark reproduces the qualitative pattern of
\citet{cieslak2015}---a null $R^{2}$ band that widens with regressor
persistence---but not the exact published band
(\Cref{sec:meth-ehmc}). The claim that the observed in-sample $R^{2}$
``exceeds the EH null'' therefore rests on our own calibration. The margin
is wide (the realised $R^{2}$ sit far above the simulated 95th percentile),
which is why we regard the conclusion as safe, but the gap is documented
rather than closed.

\paragraph{Conditionality of the real-time result.} Finally, the robustness
chapter itself draws the sharpest boundaries. The global factor's
out-of-sample edge is a property of \emph{expanding-window} estimation and
disappears under ten- and fifteen-year rolling windows
(\Cref{sec:rob-oos-scheme}); it depends materially on measuring the trend
with \emph{core} rather than headline inflation
(\Cref{sec:rob-core-headline}); and it fails inside the 2008 crisis window.
The FX-adjusted factor's thin pooled gain is concentrated post-2009 and is
sensitive to the training minimum. These are not footnotes: they delimit
the population of investors---patient, long-memory, core-inflation-based
estimators of a slow factor---for whom the real-time predictability exists.

\section{Implications}
\label{sec:disc-implications}

\paragraph{For the international bond investor.} The portfolio prescription
that follows from the evidence has three parts. First, \emph{hedge the
currency}. The predictable component of international bond returns lives in
local-currency space; dollar conversion destroys roughly two-thirds of it
in sample and almost all of it out of sample, and the timing strategy's
Sharpe ratio falls from $0.38$ hedged to about $0.25$ unhedged
(\Cref{ch:strategy}). Currency exposure in a G10 bond portfolio is, from
the standpoint of cycle-based predictability, uncompensated noise. Second,
\emph{time the aggregate, not the cross-section}. The factor that survives
real time is a single global series; it supports varying overall duration
exposure---which raised the hedged Sharpe ratio from $0.31$ to $0.38$ with
a certainty-equivalent gain of about half a percent a year, achieved mainly
by sidestepping drawdowns such as 2022---but the out-of-sample evidence
gives no licence for country selection on local cycle factors. Third,
\emph{respect the factor's limits}: it is slow, it requires a long
estimation history, and it carries no protection in a flight-to-quality
panic, when its real-time forecasts were at their worst.

\paragraph{For risk management and the periphery.} The Italian result has a
practical reading for holders of peripheral euro-area debt: under sovereign
stress, local cycle deviations regain forecasting power that the global
factor misses entirely, so a risk manager monitoring peripheral exposure
should track the \emph{local} factor exactly when spreads widen---the one
regime in which it demonstrably worked in real time. Integration-based risk
models that price peripheral bonds off a common factor will be most wrong
precisely when the position is most dangerous.

\paragraph{For the term-premium literature.} The findings carry two
messages back to the research literature. First, they add G10 evidence that
term premia are anchored by slow macro trends---the international
counterpart of the falling-stars mechanism of \citet{bauer2020}, and a
factor-level complement to the trend-augmented forecasting of
\citet{zhang2021}: it is the \emph{deviation} from the trend, aggregated
globally, that is priced, and it is priced in common across markets. Second,
they suggest that the in-sample horse race between forward-rate and
macro-anchored predictors is the wrong contest. In sample the forward
factor wins; in real time only the macro-anchored global factor survives.
Studies that evaluate candidate state variables purely in sample will
systematically favour price-based over macro-anchored predictors, and the
out-of-sample discipline of \citet{campbellthompson2008}---applied here to
the full generated-regressor chain, not only to the final
regression---reverses that verdict. A natural next step, beyond the scope
of this thesis, is to embed the global cycle in a no-arbitrage term
structure model with a shared international trend, closing the gap between
the reduced-form evidence assembled here and the structural falling-stars
framework.
